\documentclass[11pt]{article}
\usepackage[margin=1in]{geometry}
\usepackage{hyperref}
\hypersetup{colorlinks=true,linkcolor=black,urlcolor=black}
\pagestyle{empty}
\setlength{\parindent}{0pt}
\setlength{\parskip}{0.9em}

\begin{document}

Dear Editors,

We are pleased to submit our manuscript entitled ``Competing feedback channels drive phase transitions in collective emotion'' for consideration as a Research Article in \textit{Nature Computational Science}.

A central open problem in computational social science is predicting when collective sentiment will destabilize---shifting from prolonged calm to self-sustaining polarization---and what structural conditions govern this transition. Existing computational approaches to online polarization, including agent-based opinion models and data-driven cascade analyses, typically assume a single information channel. Real digital platforms, however, operate as mixed ecologies in which stabilizing mainstream media compete with amplifying user-generated content, creating opposing feedback forces whose balance can shift over time.

We develop a computational framework that formalizes this channel-balance competition. The framework integrates three levels of analysis: (1)~mean-field theory that yields an analytic critical point and a Ginzburg--Landau effective potential, providing closed-form predictions for when the neutral state loses stability; (2)~agent-based simulations on heterogeneous networks that confirm the transition structure persists under finite-size fluctuations and topological heterogeneity; and (3)~empirical validation in large-scale social-media data from the COVID-19 period, where both predicted signatures---activity-driven polarization jumps and media-composition-driven volatility---appear. A direction--intensity decomposition offers a computable early-warning indicator: elevated engagement activity signals approaching instability before polarization emerges.

We believe this work aligns with \textit{Nature Computational Science}'s commitment to computational social science. The contribution is methodological---a new analytical and computational toolkit for understanding collective transitions in social systems---grounded in empirical data and offering actionable diagnostics for platform governance. All simulation code and aggregated data will be publicly available upon acceptance, consistent with the journal's commitment to reproducibility.

We confirm that this manuscript has not been published elsewhere and is not under consideration by another journal. All authors have approved the submission.

Thank you for considering our work.

Sincerely,

\textit{On behalf of all authors,}

\vspace{0.5em}

\textbf{Corresponding authors:}

Jinlin Wu\\
Urban Governance and Design Thrust\\
The Hong Kong University of Science and Technology (Guangzhou)\\
Email: jwu923@connect.hkust-gz.edu.cn

Zhihang Liu\\
Institute of Space and Earth Information Science\\
The Chinese University of Hong Kong\\
Email: zhihangliu@cuhk.edu.hk

\end{document}
